\section{Evaluation}
\label{sec:evaluation}
We test quality at fixed physical memory, complete request latency, and quality
preservation at smaller measured points. All MosaicKV results execute
\mexact; prototypes, residual storage, and repair are disabled. Baselines
marked R are local reimplementations.

\subsection{Setup, controls, and metrics}
\begin{table*}[t]
\caption{Completed evaluation strata. Native-retention and shared-byte lanes
are separate. Revisions are abbreviated; full identifiers are in the artifact.
Each GPU worker uses one A100-SXM4-80GB and 12 CPU cores.}
\label{tab:completed-setup}
\centering\small
\begin{tabularx}{\textwidth}{P{2.2cm}P{3.3cm}P{3.2cm}Y}
\toprule
Lane / cohort & Model revision & Measured operating points & Samples and metrics\\\midrule
COCO & LLaVA-v1.6-Mistral-7B, \texttt{2424fdd4} & 10/20/40/80 MiB, FP16 & 256 images; corpus CIDEr\\
MileBench E6 & LLaVA-1.5-7B, \texttt{b234b804} & 16/20/28/40/80 MiB, FP16 & 192 total; 48 per subtask; Accuracy / ROUGE-L\\
MileBench earlier & Same LLaVA-1.5 revision & 10/20/40/80 MiB, FP16 & 256 total; 64 per subtask; separate appendix tables\\
PrefixKV native & Same LLaVA-1.5 revision & Nominal retention 0.5, BF16 & 1,000 descriptions; 218 paper-style MM-Vet; PPL / reference ROUGE-L\\
PrefixKV shared & Same LLaVA-1.5 revision & Exactly 40 MiB, FP16 & Same 1,218 inputs; precision-matched FullKV reference\\
\bottomrule
\end{tabularx}
\end{table*}

\paragraph{Inputs and execution.}
COCO Caption uses revision \texttt{c3213b54}; MileBench uses
\texttt{53c7a580} \citep{lin2014coco,song2024milebench}. MileBench includes
ActionPrediction, ActionSequence, CLEVR-Change, and Spot-the-Diff: two temporal
Accuracy subtasks and two semantic ROUGE-L subtasks. This is a representative
selection, not the complete benchmark. Multi-image inputs use frozen
materialized preprocessing, including merged-image transport where required;
all methods receive identical processed content. PrefixKV uses its original
description and 218-question routes, not judge-scored MM-Vet
\citep{wang2025prefixkv,yu2023mmvet}.

Within each shared-byte lane, checkpoint/revision, processor/chat template,
prompts, media, tokenization, deterministic greedy generation, precision, eager
attention, one-token blocks, mandatory tokens, and output length are fixed.
COCO/MileBench generate 64 tokens; PrefixKV routes generate 1,024.
COCO protects the last prompt token; MileBench and shared PrefixKV use the
first/last policy. Native PrefixKV preserves its method-specific maintenance,
so its native comparison is not byte-matched. Runs use the pinned common
environment (PyTorch 2.8.0, Transformers 4.55.2) and five warmups.
Times are means over requests with one measured quality trial, not repeated
microbenchmarks or production throughput.

\paragraph{Physical memory.}
Shared-byte analysis verifies equal initial, final, and peak live KV storage
per sample and the largest feasible allocation below each cap. Configured caps
alone are insufficient. COCO decoding adds 7.875 MiB; MileBench adds 31.5 MiB;
the PrefixKV append-only shared lane adds 511.5 MiB. Thus the same 40 MiB
initial cap yields 47.875, 71.5, or 551.5 MiB after generation.
FullKV supplies identity and quality references and is excluded from latency
comparisons.

\paragraph{Metrics and uncertainty.}
COCO uses the pinned lmms-eval PTB tokenizer and CIDEr implementation
\citep{vedantam2015cider}. Corpus CIDEr is recomputed over the complete cohort
and inside every one of 10,000 paired bootstrap replicates; shard scores are
never averaged. MileBench recomputes its equal-dataset category hierarchy in
every replicate, keeping Accuracy and ROUGE-L separate. PrefixKV reports
\begin{equation}
 {\rm PPL}=\exp\!\left(\frac{\sum_s{\rm NLL}_s}{\sum_s T_s}\right)
\end{equation}
and mean FullKV-reference ROUGE-L F1 \citep{lin2004rouge}. PPL covers all
138,905 description tokens and 20,949 question-answer tokens. Teacher-forced
annotation targets differ from free-running generation outputs.
We report paired 95\% intervals \citep{efron1979bootstrap} and family-wise
adjusted intervals for the registered MileBench lower-memory tests.
No average mixes different metric units.

\paragraph{Completeness and evidence status.}
COCO contains 3,328 complete rows, MileBench E6 3,072, the earlier MileBench
cohort 3,328, and PrefixKV 8,526 across seven arms. Numerical exports contain
no failed or duplicate rows. COCO and the earlier MileBench pool have prior
protocol use and are not pristine held-out evidence. E6 uses an
exclusion-reviewed new cohort relative to its bound ledgers; cohorts are
reported separately. Numerical export independently checks completeness and
per-example shared storage; artifact details appear in the separate appendix.

\subsection{COCO: shared-byte quality and runtime}
\begin{table}[t]
\caption{COCO at 40 MiB. Storage is MiB; times are seconds. FullKV is a
quality reference. R denotes a local reimplementation.}
\label{tab:coco-measured}
\centering\small\resizebox{\columnwidth}{!}{% Generated from hash-verified measured artifacts by scripts/export_results.py.
\begin{tabular}{lrrrrr}
\toprule
Method & CIDEr $\uparrow$ & Init. & Final & Decode & E2E\\\midrule
FullKV & 111.24 & 285.61 & 293.481 & --- & ---\\
Prompt Top-K & 111.59 & 40.00 & 47.875 & 2.95 & 17.22\\
VL-Cache (R) & 107.84 & 40.00 & 47.875 & 3.12 & 20.75\\
MosaicKV exact & 111.48 & 40.00 & 47.875 & 2.95 & 16.68\\
\bottomrule
\end{tabular}
}
\end{table}

At identical 40 MiB initial and 47.875 MiB final/peak cache,
MosaicKV scores 111.48 CIDEr versus 107.84 for VL-Cache (R), a 3.64-point
gain with paired interval $[0.53,7.06]$. Mean request time falls from
20.75 to 16.68 s (19.65\%); the paired absolute reduction is
$[3.92,4.24]$ s. Decode falls from 3.12 to 2.95 s (5.53\%).
Most of the 4.08 s request-time difference occurs outside decode.

Prompt Top-K is a stronger quality competitor: it scores 111.59.
MosaicKV's difference is $-0.11$ CIDEr, interval $[-0.63,0.31]$, with
3.14\% lower mean request time. Quality ordering is unresolved.
Against FullKV, the difference is $+0.24$, interval $[-3.74,4.51]$.
The cache is 86.0\% smaller initially and 83.7\% smaller after decoding than
FullKV, but those point estimates do not establish a zero-loss quality guarantee.

\begin{figure*}[t]
\centering\begin{tikzpicture}
\begin{groupplot}[group style={group size=3 by 1,horizontal sep=13mm},
 width=0.31\textwidth,height=4.7cm,
 xlabel={Initial active KV (MiB)},grid=major,grid style={gray!18},
 tick label style={font=\scriptsize},label style={font=\small},
 title style={font=\small},legend style={font=\scriptsize},
 legend columns=3]
\nextgroupplot[title={COCO: 256 images},ylabel={CIDEr $\uparrow$},xtick={10,20,40,80},ymin=88,ymax=114,legend to name=qualitylegend]
\addplot[exactblue,thick,mark=*] table[x=cap,y=score]{generated/coco_mosaickv_exact.dat};
\addlegendentry{MosaicKV exact}
\addplot[black,dashed,mark=square*] table[x=cap,y=score]{generated/coco_prompt_attention_topk.dat};
\addlegendentry{Prompt Top-K}
\addplot[protogold,thick,mark=triangle*] table[x=cap,y=score]{generated/coco_vl_cache_reimpl.dat};
\addlegendentry{VL-Cache / LOOK-M (R)}
\nextgroupplot[title={MileBench E6: semantic},ylabel={ROUGE-L $\times100$ $\uparrow$},xtick={16,40,80},ymin=11,ymax=17]
\addplot[exactblue,thick,mark=*] table[x=cap,y=rouge]{generated/mile192_mosaickv_exact.dat};
\addplot[black,dashed,mark=square*] table[x=cap,y=rouge]{generated/mile192_prompt_attention_topk.dat};
\addplot[protogold,thick,mark=triangle*] table[x=cap,y=rouge]{generated/mile192_lookm_reimpl.dat};
\nextgroupplot[title={MileBench E6: temporal},ylabel={Accuracy (\%) $\uparrow$},xtick={16,40,80},ymin=0,ymax=29]
\addplot[exactblue,thick,mark=*] table[x=cap,y=accuracy]{generated/mile192_mosaickv_exact.dat};
\addplot[black,dashed,mark=square*] table[x=cap,y=accuracy]{generated/mile192_prompt_attention_topk.dat};
\addplot[protogold,thick,mark=triangle*] table[x=cap,y=accuracy]{generated/mile192_lookm_reimpl.dat};
\end{groupplot}
\node[anchor=north] at ($(group c1r1.south)!0.5!(group c3r1.south)+(0,-10mm)$) {\pgfplotslegendfromname{qualitylegend}};
\end{tikzpicture}

\caption{Measured fixed-budget curves. Left: COCO, 256 images; center/right:
the separate 192-example MileBench E6 cohort. Every compressed method runs at
every displayed cap. Lines connect measured points, without establishing
quality between them. ROUGE-L is scaled by 100; Accuracy and CIDEr retain
their respective meanings.}
\label{fig:curves}
\end{figure*}

Figure~\ref{fig:curves} retains every tested point. At 10 MiB, MosaicKV
improves CIDEr over VL-Cache (R) by 11.34, interval $[6.55,16.49]$;
at 20 MiB the gain is 4.35, interval $[0.01,8.90]$. At 80 MiB the
0.75-point gain is unresolved. Prompt Top-K tracks MosaicKV closely.
The result supports advantages over this VL-Cache port under tight capacity,
rather than universal dominance over selectors.

\subsection{MileBench: complementary evidence at fixed memory}
At 16 MiB, MosaicKV reaches 21.88\% Accuracy and 14.52 ROUGE-L points,
versus LOOK-M (R)'s 3.12\% and 12.31. Paired gains are $+18.75$ Accuracy
points, interval $[10.42,27.08]$, and $+2.21$ ROUGE-L points,
interval $[0.80,3.62]$. Mean request time falls from 24.36 to 14.95 s
(38.6\%). This tight-budget collapse concerns the byte-capped
reimplementation regime, not official LOOK-M's original paper setting.

\begin{table}[t]
\caption{MileBench E6 at 40 MiB initial / 71.5 MiB final and peak cache.
ROUGE-L is scaled by 100. Each metric covers its 96 applicable examples.}
\label{tab:mile40}
\centering\small\resizebox{\columnwidth}{!}{% Generated from hash-verified measured artifacts by scripts/export_results.py.
\begin{tabular}{lrrr}
\toprule
Method & Acc. (\%) $\uparrow$ & R-L $\uparrow$ & E2E s $\downarrow$\\\midrule
Prompt Top-K & 22.92 & 15.88 & 18.03\\
LOOK-M (R) & 25.00 & 14.48 & 34.22\\
MosaicKV exact & 22.92 & 15.85 & 16.71\\
\bottomrule
\end{tabular}
}
\end{table}

At 40 MiB, MosaicKV improves ROUGE-L over LOOK-M (R) by 1.36 points,
interval $[0.15,2.55]$, while Accuracy is 2.08 points lower,
interval $[-7.29,3.12]$. Its Accuracy matches Prompt Top-K on paired
correctness; its ROUGE-L difference is $-0.04$, interval $[-0.25,0.17]$.
Mean time is 51.2\% lower than LOOK-M (R) and 7.33\% lower than Top-K.
Across E6 caps, MosaicKV is 38.6--58.4\% faster than LOOK-M (R) and
5.6--7.5\% faster than Top-K.

The earlier 256-example cohort measured a similar trade-off at 40 MiB:
equal Accuracy point estimates against LOOK-M (R), $+0.98$ ROUGE-L points
with interval $[0.01,2.01]$, and 50.9\% lower mean time.
At 80 MiB, however, LOOK-M (R)'s Accuracy point estimate was higher.
The appendix preserves both cohorts' complete and unfavorable points;
they are not pooled.

\subsection{PrefixKV: native fidelity and shared capacity}
\begin{table*}[t]
\caption{Completed PrefixKV native-retention lane, BF16, registered nominal
retention 0.5. R-L is FullKV-reference ROUGE-L scaled by 100; PPL is corpus
perplexity. Storage is MiB. Unequal realized bytes and decode maintenance make
this a native operating-point comparison, not an equal-budget policy comparison.}
\label{tab:prefixnative}
\centering\small% Generated from hash-verified measured artifacts by scripts/export_results.py.
\begin{tabular}{llrrrr}
\toprule
Task & Method & PPL $\downarrow$ & R-L $\uparrow$ & Init. & Peak\\\midrule
LLaVA-Desc. & PrefixKV (R) & 3.0889 & 63.42 & 149.7 & 406.4\\
LLaVA-Desc. & MosaicKV exact & 2.9620 & 71.09 & 102.6 & 614.1\\
MM-Vet & PrefixKV (R) & 5.3339 & 64.97 & 151.0 & 407.7\\
MM-Vet & MosaicKV exact & 5.2775 & 79.03 & 103.6 & 615.1\\
\bottomrule
\end{tabular}

\end{table*}

The four-example official-path gate passed position, per-layer-length,
retained-byte, token, NLL, and ROUGE-L parity. The large comparison uses the
pinned profile through a local port; it is not an official 1,218-example
execution. Native PrefixKV retains ongoing eviction.

On LLaVA-Description, MosaicKV improves PPL from 3.0889 to 2.9620 and
ROUGE-L from 63.42 to 71.09. The paired PPL difference is $-0.1269$,
interval $[-0.1328,-0.1211]$; the ROUGE-L gain is 7.67 points,
interval $[6.62,8.70]$. On MM-Vet, ROUGE-L improves from 64.97 to 79.03
(gain 14.04, interval $[11.22,16.76]$). PPL improves from 5.3339 to
5.2775 in point estimate, but its paired interval crosses zero.

MosaicKV uses about 31.4\% less initial KV at these native points.
The saving reverses during generation: its append-only peak is
614--615 MiB; PrefixKV's ongoing eviction limits its peak to 406--408 MiB.
MosaicKV's mean request time is 20.7--20.9\% higher. Thus this comparison
supports stronger fidelity at the registered point, not lower final memory
or latency.

In the separate FP16 shared 40 MiB lane, compressed methods finish at
551.5 MiB. MosaicKV and Top-K have close quality: description ROUGE-L is
63.09 versus 63.20; question ROUGE-L is 72.71 versus 72.62.
Description PPL is slightly but significantly worse for MosaicKV
($+0.0015$); question PPL differences are unresolved.
MosaicKV lowers mean time by 0.75--0.89\%.
The byte-rescaled PrefixKV adaptation collapses to roughly 14 ROUGE-L points.
The appendix reports it as an adaptation stress test; it does not characterize
native PrefixKV.

\subsection{Lower-memory quality equivalence: measured limits}
\begin{figure}[t]
\centering\begin{tikzpicture}
\begin{axis}[width=\columnwidth,height=4.5cm,xmin=.8,xmax=3.2,
 xtick={1,2,3},xticklabels={Initial,Final,Peak},ylabel={Live KV (MiB)},
 tick label style={font=\scriptsize},label style={font=\small},grid=major,
 grid style={gray!18},legend style={font=\scriptsize,at={(.02,.98)},anchor=north west}]
\addplot[exactblue,thick,mark=*] table[x=boundary,y=mosaickv_exact]{generated/prefix_boundaries.dat};
\addlegendentry{MosaicKV exact}
\addplot[protogold,thick,mark=triangle*] table[x=boundary,y=prefixkv_reimpl]{generated/prefix_boundaries.dat};
\addlegendentry{PrefixKV native (R)}
\end{axis}
\end{tikzpicture}

\caption{Native PrefixKV lane on LLaVA-Description: memory ordering reverses
after long decoding. Lines connect accounting boundaries, not time samples.
Smaller initial KV is insufficient to claim smaller final KV.}
\label{fig:boundaries}
\end{figure}

The E6 protocol chose cross-budget pairs using storage before opening quality.
MosaicKV 16 versus baseline 40 MiB saves 60\% initially and 33.57\%
at final/peak; 40 versus 80 MiB saves 50\% initially and 35.87\%
at final/peak. All tested pairs pass memory but fail zero-margin quality
non-inferiority after family-wise correction. Against LOOK-M (R) 80 MiB,
MosaicKV 40 MiB has positive point differences on both metrics, yet adjusted
intervals cross zero. Quality preservation remains unresolved.

COCO has an exploratory point-estimate example: MosaicKV 20 MiB scores
107.97 versus VL-Cache (R) 40 MiB at 107.84, with 41.8\% lower final KV.
The artifact has no registered cross-budget confidence test proving
equivalence, and Top-K also reaches that quality at 20 MiB.
The larger VL-Cache point is therefore not the nearest quality-eligible
competitor. Figure~\ref{fig:boundaries} shows another limit: maintenance
policy can outweigh initial savings. All registered outcomes appear in
the appendix.

\subsection{Why these operating regimes favor MosaicKV}
\begin{table}[t]
\caption{MileBench E6 phase means at 40 MiB. Compression includes recorded
planning and construction work. Other request phases account for the
difference between compression plus decode and end-to-end time.}
\label{tab:phases}
\centering\small\resizebox{\columnwidth}{!}{% Generated from hash-verified measured artifacts by scripts/export_results.py.
\begin{tabular}{lrrr}
\toprule
Method & Compress s & Decode s & E2E s\\\midrule
Prompt Top-K & 11.63 & 2.79 & 18.03\\
LOOK-M (R) & 26.36 & 2.80 & 34.22\\
MosaicKV exact & 10.50 & 2.77 & 16.71\\
\bottomrule
\end{tabular}
}
\end{table}
At equal bytes, quality differences reflect retained evidence and
transformations, not extra storage. Prompt queries emphasize current demand;
graph coverage discourages redundant neighborhoods; value signals distinguish
information beyond attention mass. Exact retention avoids errors from unsafe
merging. These mechanisms are consistent with COCO gains over VL-Cache (R)
and tight-budget MileBench gains over LOOK-M (R). The near tie with Top-K
also shows how much useful signal prompt attention already captures.

The execution explanation is more direct: most heads bypass graph work;
one-token selection, sparse edge assembly, and packing avoid full-source
Python planning. Figure~\ref{fig:latency} includes complete compression costs.
The VL-Cache port's absent fused kernels and LOOK-M port's merge/construction
work also affect runtime. At equal bytes, decode gains cannot be credited to
reading fewer KV bytes; layout and implementation differ. These are
mechanistic interpretations, not causal ablations of each component.
Table~\ref{tab:phases} makes the bottleneck explicit: MosaicKV compression is
10.50 s versus 11.63 s for Top-K and 26.36 s for LOOK-M (R), while their
decode means are close (2.77--2.80 s). Thus most of the measured advantage
comes from planning/construction and other pre-decode work. Compression still
dominates MosaicKV itself; reducing it remains an important systems target.

\begin{figure}[t]
\centering\begin{tikzpicture}
\begin{axis}[width=\columnwidth,height=4.5cm,xlabel={Initial active KV (MiB)},
 ylabel={Mean request time (s) $\downarrow$},xtick={16,20,40,80},
 tick label style={font=\scriptsize},label style={font=\small},grid=major,
 grid style={gray!18},legend style={font=\scriptsize,at={(.02,.98)},anchor=north west}]
\addplot[exactblue,thick,mark=*] table[x=cap,y=latency]{generated/mile192_mosaickv_exact.dat};
\addlegendentry{MosaicKV exact}
\addplot[black,dashed,mark=square*] table[x=cap,y=latency]{generated/mile192_prompt_attention_topk.dat};
\addlegendentry{Prompt Top-K}
\addplot[protogold,thick,mark=triangle*] table[x=cap,y=latency]{generated/mile192_lookm_reimpl.dat};
\addlegendentry{LOOK-M (R)}
\end{axis}
\end{tikzpicture}

\caption{MileBench E6 mean request latency at identical physical cache sizes,
including compression. Request-level observations use one quality trial per
sample and are not throughput measurements.}
\label{fig:latency}
\end{figure}
